\documentclass[11pt,a4paper]{article}
\usepackage[T1]{fontenc}
\usepackage[utf8]{inputenc}
\usepackage{lmodern,amsmath,amssymb,amsthm,mathtools,microtype,mathrsfs}
\usepackage[margin=25mm,top=23mm,bottom=25mm]{geometry}
\usepackage{booktabs,array,longtable,enumitem,xcolor,fancyhdr}
\usepackage[colorlinks=true,linkcolor=blue!45!black,urlcolor=blue!45!black,citecolor=blue!45!black]{hyperref}
\hypersetup{pdftitle={Reading the Alpoge--Claude RH and Dirichlet-GRH proof: working version 0.31},pdfauthor={ChatGPT},pdfsubject={Partial human-readable reconstruction with explicit arithmetic inputs},pdfkeywords={Riemann hypothesis, Lean, theta lifts, prime sums, Arakelov slopes}}
\setlength{\parindent}{0pt}
\setlength{\parskip}{5pt plus 1pt}
\setlist{topsep=3pt,itemsep=3pt,leftmargin=*}
\setlength{\headheight}{14pt}
\pagestyle{fancy}\fancyhf{}
\fancyhead[L]{\small RH and Dirichlet GRH: a reconstruction}
\fancyhead[R]{\small Working version 0.31}
\fancyfoot[C]{\thepage}
\renewcommand{\headrulewidth}{0.3pt}
\numberwithin{equation}{section}
\newtheorem{theorem}{Theorem}[section]
\newtheorem{proposition}[theorem]{Proposition}
\newtheorem{lemma}[theorem]{Lemma}
\newtheorem{corollary}[theorem]{Corollary}
\theoremstyle{remark}\newtheorem{remark}[theorem]{Remark}
\newcommand{\Rea}{\operatorname{Re}}
\newcommand{\HS}{\mathrm{HS}}
\newcommand{\tr}{\operatorname{tr}}
\newcommand{\codim}{\operatorname{codim}}
\newcommand{\covol}{\operatorname{covol}}
\newcommand{\OK}{\mathcal O_K}
\newcommand{\src}[2]{\href{https://github.com/CocoSgt/riemann-hypothesis-lean/blob/7cbf7a0c32fca1badf05931879a37c3f8a2b03b4/src/#1}{\texttt{#2}}}
\newcommand{\status}[1]{\textbf{Status:} #1}
\begin{document}
\section{Compatible design choices and the remaining existence obligation}
\label{sec:design-existence-v031}

\status{DEX1--DEX2 install the literal family thinning condition from the
already reconstructed centered-character estimates. DEX3 matches the
published Minkowski theorem, with its exact integer, and DEX4 proves the
required logarithmic bound. DEX5 proves the numerical proviso directly
from the original rank formula. DEX6--DEX11 reconstruct the common-threshold,
common-selector and actual-family joins. The final existence statement
remains conditional on the specified original recursive objects, rank
formula and level-prime size supplier. None of these joins asserts that
independently existing geometric or theta objects are compatible.}

Throughout, a family means the original \texttt{FamP}, with its actual
fields, maps, word labels, lattices and metrics. All statements concern the
pinned source and the fixed target character used in the global audit.
The standard background of the preceding chapters is retained. In
particular, CC10--CC18 are used in their proved scopes, including the
classical zero-free-region input already recorded there.

\paragraph{DEX1: the conductor and odd-part bounds on one actual family.}
\label{claim:DEX1}
Let $D$ be an original datum and let $\chi$ satisfy
\texttt{IsCharED}. Thus $\chi$ is primitive modulo $N$ and, on
integers prime to $M=Q_Sqd$, agrees with a character of the original
finite Artin group, pulled back by the actual Artin map. Then
\[
                         N\mid M.                         \tag{DEX.1}
\]
To see the precise conductor assertion, inflate both characters to
$\operatorname{lcm}(N,M)$. Every unit there has an integer representative
prime to both moduli, so the two inflations agree. The right-hand
character factors through $M$. By the prime-by-prime definition of the
conductor, each local conductor exponent of the primitive $\chi$ is at
most the corresponding exponent of $M$. This proves (DEX.1), including
$N=1$. Equivalently, the local kernel killed by reduction to $M$ is
killed by $\chi$; Chinese remainders separate these local assertions.

For an original family $F$, the late members satisfy
\[
 o\le C_o,\qquad Q_S\le C_o n,\qquad
 n\le o2^e,\qquad 2^e\mid Q_S .                           \tag{DEX.2}
\]
The first two are literal family fields and the last two are the
original frame and host fields. Write $Q_S=2^vQ_{S,\mathrm{odd}}$.
Since $v\ge e$, (DEX.2) gives
\[
 2^eQ_{S,\mathrm{odd}}\le Q_S
       \le C_on\le C_o^2 2^e,
 \qquad Q_{S,\mathrm{odd}}\le C_o^2.                      \tag{DEX.3}
\]
Here $C_o\ge0$ follows from $o\le C_o$ and $o\ge0$.
In particular $B=\lceil C_o^2\rceil$ bounds the odd part uniformly
over the late members. This does not bound the power of two in $Q_S$.

Finally, \texttt{BoxE} at any supplied scale $Y$ implies the following
real-zero exclusion: every real primitive host character whose conductor
is divisible by $q$ or $d$ has no real zero $\beta>1-1/400$.
Indeed, if $d\mid N$, use the first \texttt{GoodFor} clause. Otherwise
$q\mid N$ and the second clause applies. In both cases the height of the
real argument is zero, which is at most the nonnegative box height.
No comparison between $Y$ and the large variable of the character sum is
needed for this real restriction.

\paragraph{DEX2: close \texttt{ChosenThin}, with coefficient one.}
\label{claim:DEX2}
For every original family $F$ and $\epsilon>0$, we prove
\[
 \mathrm{BoxesF}(F)\ \wedge\ \mathrm{QDRowsF}(F)
               \quad\Longrightarrow\quad
                         \mathrm{ChosenThin}(F,1).       \tag{DEX.4}
\]
In particular the extra \texttt{ChosenQD} premise appearing in one
source wrapper is unnecessary for this implication.

Fix $F$. Beyond one threshold, \texttt{BoxesF} supplies a scale $Y'$
and \texttt{BoxE} for each actual $D=F.D(x,i)$. The other hypothesis
supplies a possibly different scale $Y$, with
$x^{1/J}\le Y\le x$, and its conductor bound
\[
       \log(Q_Sqd)\le2Y^{\epsilon/5}\log^2Y.
\]
The family has $J\ge1$. After imposing $x>1$, we have $Y>1$;
both $Y^{\epsilon/5}$ and $\log^2Y$ increase on this range. Hence
\[
                     \log(Q_Sqd)
                         \le2x^{\epsilon/5}\log^2x.      \tag{DEX.5}
\]
The two scales need not be equal: the box is used only through its
real-zero exclusion, and the conductor inequality only through (DEX.5).

Take $B$ from DEX1. CC18 supplies $c_s>0$ excluding real zeros in
$(1-c_s,1)$ for every nonprincipal character of conductor at most $8B$.
Set $c_0=\min(c_s,1/400)>0$. For any real primitive nonprincipal
\texttt{IsCharED} character, there are two cases. If $q\mid N$ or
$d\mid N$, DEX1 gives the required exclusion. Otherwise (DEX.1) and
the primality of $q,d$ give $N\mid Q_S$. CC17 gives two-adic conductor
exponent at most three, so $N\mid8Q_{S,\mathrm{odd}}$ and $N\le8B$.
The finite-conductor exclusion then applies. Thus every such character
satisfies the precise \texttt{NoExc} premise of CC16 with this common
$c_0$. Nonreal and primitive principal characters require no additional
real-zero exclusion in that premise.

By (DEX.1) and (DEX.5), every character also has
$\log N\le2x^{\epsilon/5}\log^2x$. Apply CC16 with $A=2$ and $c'=1$.
CC14 identifies its centered function with the literal
\texttt{GlobalTheoryTheory.gChi}. We obtain, simultaneously for all
late $x$, all $i$, all primitive host characters, and every
$y\ge e^{x^{\epsilon/3}}/3$,
\[
             |g_\chi(y)|\le y\exp(-x^{\epsilon/8}).        \tag{DEX.6}
\]
The threshold may depend on $F$ through $C_o$, as the definition of
\texttt{ChosenThin} permits. The coefficient $1$ does not. Taking the
maximum with $F.x_1+1$ and the two supplied family thresholds proves
(DEX.4) with the exact universal quantifiers. A family itself supplies
$\epsilon>0$ through any index at $x=0$, so the pointwise implication
also yields \texttt{Closed} on any class: choose outer bound $1$ and
the zero transported threshold. Consequently no separate literal
\texttt{LitThin} package is needed at this consumer; the analytic inputs
of CC retain their previously stated roles.

\paragraph{DEX3: the exact published Minkowski application.}
\label{claim:DEX3}
For $n\ge0$ put
\[
 a_p(n)=\sum_{j\ge0}\left\lfloor\frac{n}{p^j(p-1)}\right\rfloor,
 \qquad M(n)=\prod_{p\le n+1}p^{a_p(n)}.                  \tag{DEX.7}
\]
These are precisely the source's \texttt{minkowskiExp} and
\texttt{minkM}. Terms with $j\ge n+1$ vanish: $p^j\ge2^j>n$.
Primes larger than $n+1$ contribute zero. Thus the finite truncations
in the source do not change either expression.

The published input is Minkowski's theorem in Serre,
\emph{Bounds for the orders of the finite subgroups of $G(k)$},
Lecture I, \S1.1, Theorem 1(i): for a finite subgroup
$H\subset\mathrm{GL}_n(\mathbb Q)$, the $p$-adic valuation of $|H|$
is at most $a_p(n)$. See
\url{https://www.college-de-france.fr/media/jean-pierre-serre/UPL3821667391778701726_6___Bounds_for_the_orders.pdf},
printed p.~4. The source \texttt{MinkMult} has exactly this rational
general linear group and assumes the subgroup finite. Applying the
theorem at every prime gives $|H|\mid M(n)$. For $n=0$ the group is
trivial and $M(0)=1$. This checks every rank, including the case not
covered by the published statement $n\ge1$, and proves the literal
\texttt{MinkDiv}. No estimate over a larger coefficient field is
substituted for this rational theorem.

\paragraph{DEX4: the exact denominator-size constant twelve.}
\label{claim:DEX4}
The numerical consequence of (DEX.7) has an independent elementary
proof. Legendre's formula and $p\le2(p-1)$ give, term by term,
\[
 \left\lfloor\frac{n}{p^j(p-1)}\right\rfloor
 \le \left\lfloor\frac{2n}{p^{j+1}}\right\rfloor,
 \qquad M(n)\mid(2n)! .                                 \tag{DEX.8}
\]
Legendre's formula itself counts, among $1,\ldots,2n$, one contribution
for each multiple of $p$, a second for each multiple of $p^2$, and so
on. Since the prime-power factors are coprime for distinct primes,
the valuation inequalities prove the displayed divisibility. For
$n\ge1$, it follows that
\[
                       \log M(n)\le2n\log(2n).
\]
Write $z=fN(m)$ for the actual rank product. The original ranks have
$z\ge2$. Therefore
\[
 \log M(2z)\le4z\log(4z)
       =4z(2\log2+\log z)\le12z\log z.                 \tag{DEX.9}
\]
If \texttt{DGamDvdF} holds, then eventually every original positive
integer $D_\Gamma(u)$ divides this same $M(2fN(m))$. Divisibility by
a positive integer ensures $D_\Gamma(u)>0$ and bounds it above by that
integer. Taking logarithms in (DEX.9) proves \texttt{DGamSize} with
the literal constant $12$, uniformly over all words and family indices.
The original exceptional-prime support is unchanged: all prime divisors
of $M(n)$ are at most $n+1$ by (DEX.7). This is a consequence of the
supplied divisibility, not a new construction of the original lattice.

\paragraph{DEX5: the numerical proviso from the actual rank formula.}
\label{claim:DEX5}
Let $D,L$ be original data and assume the source \texttt{RankM} identity
\[
 d_M(u)=d_*(u)(T_u^k-1)^f.                              \tag{DEX.10}
\]
Here $d_M(u)=D.duS(u)$, and $d_*(u)$ is the exact standard-lattice
quantity \texttt{duStarU}. Retain the original equalities
\texttt{FactDuStar} and \texttt{FactDstar}. They identify the positive
point quantities with $d_*(u)$ and $D_*$, without an isomorphism choice.
Put $\Gamma=\mathrm{GamL}(L)$ and
\[
 Y=\left(\frac{D_*}{d_*(u)}\right)^{1/(fk)}
                  \Gamma^{-2/(fk)} .                    \tag{DEX.11}
\]
The size datum's \texttt{Gam\_eq} identifies this $\Gamma$ with the
point's $\mathrm{GamP}$; thus $Y$ is the actual $T_*$.
Its original prime window says
$Y\le T_u\le(1+1/(fk))Y$, and $T_u\ge2$.
The rank fields give $1\le f\le k$ and hence
$2f\le2k\le2^k$; the last inequality follows by induction for $k\ge1$.
Using $N=fk$ and positive quantities throughout, (DEX.10)--(DEX.11)
give the exact ratio
\[
 \frac{d_M(u)\Gamma^2}{D_*}
       =\left(\frac{T_u}{Y}\right)^N(1-T_u^{-k})^f.       \tag{DEX.12}
\]
Bernoulli's inequality gives
$(1-T_u^{-k})^f\ge1-fT_u^{-k}\ge1/2$, while
$(1+1/N)^N\le e$ follows from $\log(1+t)\le t$.
Consequently
\[
                  \frac12\le\frac{d_M(u)\Gamma^2}{D_*}
                                      \le e<3.           \tag{DEX.13}
\]
For completeness, $e<3$ follows from its series:
$n!\ge2^{n-1}$ for $n\ge2$, with strict inequality at $n=3$, so
$\sum_{n\ge2}1/n!<1$. Thus (DEX.13) implies the exact
\texttt{ChProviso\_at} with $C_0=3$.

For a family, \texttt{FactDuStar} and \texttt{FactDstar} are fields of
its actual late extra record. Hence an eventual \texttt{RankM} on
that same family proves \texttt{ChosenProviso}(F,3), after taking
the maximum with $F.x_1+1$. This proof does not require an independent
formal-degree estimate. Producing (DEX.10), with its exact class-count
and coefficient-space meanings, remains a separate task.

\paragraph{DEX6: combining closed predicates before choosing the family.}
\label{claim:DEX6}
A transported lower-bound record has components
$K_{\mathrm{pre}},K_t,K_{\mathrm{mid}},K_{\mathrm{last}}$;
the last two are functions of earlier parameter lists. Define $K\vee K'$
by the pointwise maximum of corresponding components. By inspection
of the original \texttt{Large} inequalities,
\[
 F.\mathrm{Large}(K\vee K',C_H)
 \quad\Longleftrightarrow\quad
 F.\mathrm{Large}(K,C_H)\wedge F.\mathrm{Large}(K',C_H).   \tag{DEX.14}
\]
This includes the dependent arguments of the middle and last bounds:
both are evaluated on the same lists of the same $F,C_H$.

Suppose finitely many predicates are \texttt{Closed} on the same
class. Take the minimum of their positive epsilon bounds. Once
$\epsilon$ has been fixed in that range, each produces a transported
threshold; take their finite join. Any one family subsequently chosen
large for that join satisfies every predicate. Restricting the class
preserves the implication, since the original class membership is
then available. These observations prove the source operations
\texttt{Closed.and}, \texttt{Closed.sub} and implication transport.
They do not combine witnesses from different families: the uniform
closed statements precede the single existential family choice.

\paragraph{DEX7: a smaller sufficient interface for \texttt{DesignExists}.}
\label{claim:DEX7}
Fix one original class $\mathcal C$ and one function $c_T$.
The following three inputs suffice:
\[
\begin{split}
 &\mathrm{DesignChoices}(\mathcal C,3,c_T),\\
 &\mathrm{Closed}_{\mathcal C}
    [c_T(\epsilon,F.C_{YL})\le F.C_{Tu}\Rightarrow\mathrm{TuSize}(F)],\\
 &\mathrm{Closed}_{\mathcal C}
    [\exists x_0\ \forall x\ge x_0\ \forall i,
                                   \mathrm{RankM}(F.D(x,i))].
\end{split}
                                                               \tag{DEX.15}
\]
They imply the literal $\mathrm{DesignExists}(\mathcal C,3)$.
In the last line \texttt{RankM} is (DEX.10) on the original datum;
it is not a dimension match on a separately constructed model.

Use DEX5 and DEX6 to turn the last two inputs into one closed statement
giving the prime-size implication and \texttt{ChosenProviso}.
Take the minimum of its epsilon bound and that of
\texttt{DesignChoices}. For the requested $K$, first choose the
closed-statement threshold $K'$, then apply \texttt{DesignChoices}
to $K\vee K'$. It gives one $F,C_H,C_5,C_b$, with $C_H\ge0$,
$C_5,C_b>0$, the exact class membership, the required large-parameter
property, and
\[
 \begin{gathered}
 \mathrm{ThickSetEventually},\ \mathrm{ChosenTT},\
 \mathrm{ChosenBad}(C_b),\ \mathrm{ChosenQD}(C_5),\\
 \mathrm{ChosenCode}(3,C_H),\ \mathrm{BoxesF},\ \mathrm{QDRowsF},\\
 \mathrm{DGamDvdF},\ c_T(\epsilon,C_{YL})\le C_{Tu}.
 \end{gathered}
\]
The source also supplies tame distinctness, $C_F=C_H$,
\texttt{CodeUpF} and $C_{YL}\le9$; these remain attached to the same
family. DEX2 supplies \texttt{ChosenThin}(1), DEX4 supplies
\texttt{DGamSize}, and the closed statements supply
\texttt{ChosenProviso}(3) and \texttt{TuSize}.
The seven chosen clauses, with positive constants $(C_5,1,C_b)$,
are exactly $\texttt{ChosenE}(F,3,C_H)$. Equation (DEX.14) retains
$\texttt{Large}(K,C_H)$. This proves every field of
\texttt{DesignExists}, in its original order of quantifiers.

The same proof works with a supplied closed \texttt{ChosenProviso}
instead of the last line of (DEX.15). The source obtains that proviso
through \texttt{ClassNumberFormula}$\,{}_2$ and
\texttt{RankTiesF}; DEX5 records the smaller numerical input actually
used. We have not asserted that either source-specific producer exists.

\paragraph{DEX8: one selector for all eventual row conditions.}
\label{claim:DEX8}
We reconstruct the finite logical step behind \texttt{classSel}.
An eventual parameter assertion has the original order
\[
 \begin{gathered}
 \exists\epsilon_0\ \forall\epsilon,C_8,C_5,c_6,C_b\
 \exists K_0\ \forall C\\
 \exists c_1\ \forall c_D\ge c_1\
 \exists c_2\ \forall c_h\ge c_2 .
 \end{gathered}                                           \tag{DEX.16}
\]
Here the source positivity, standardness and largeness antecedents
are retained. An \texttt{EvP} assertion then has
$\kappa.\mathrm{Pos}\Rightarrow\exists x_0\forall K$ with
$\kappa.\mathrm{Of}(K)$ and $K.s.x\ge x_0$.
Two such assertions combine by a minimum for $\epsilon_0$, a join
for $K_0$, and maxima for $c_1$, then $c_2$, then $x_0$.
Each maximum is taken only after all of its permitted dependencies
have been fixed. Induction gives the finite-conjunction rule.

Write $b\in\mathrm{Bd}(D,O,c)$ for an actual bundle, and
$\mathrm{Lk}(K,\beta,c,b)$ for the source link to the original
constructed point $\beta$. The crucial distinction in
\texttt{RowE} is its universal quantifier: eventually, for
\emph{every} such linked $b$, the stated predicate $E(\kappa,D,O,c,b)$
holds. The separate \texttt{ExB} supplies one linked $b$ for each
$c$. Therefore, after the common threshold, that one $b$ satisfies
the conjunction of all finitely many row predicates. No simultaneous
choice theorem for unrelated existential bundles is being presumed.

Define $\mathrm{Built}(Q,\mathrm{Lk},\kappa,D,O,c,b)$ to mean that
there are $K,\beta,b'$ with the source context and core hypotheses,
the same link, and the literal dependent equality
\[
                 (D,O,b)=(\beta.D,\beta.O,b').            \tag{DEX.17}
\]
Let $Q_O(\kappa,D,O)$ say that, whenever $\kappa$ is positive,
for every $c$ there is such a built $b$ satisfying the combined $E$.
Take the class of original families admitting one $\kappa$ with
$\texttt{Link}(\kappa,F)$ and the source seven-duty object property
$Q_O$. At each late index choose a witness at $c=F.C_{sz}$.
This produces one selector $\sigma$. The defining witness gives
\texttt{IsSel}, including (DEX.17), and every projected row condition
holds for this same $\sigma$.

The source's additional large-base slots cause no compatibility problem.
For finitely many \texttt{RowL} conditions and one \texttt{ExBL},
replace their slots $R_j$ by $R=\bigwedge_jR_j$. The large-base
eventual statements combine by (DEX.16); at each fixed context the
finitely many eventual sets of sufficiently large bases are intersected.
Each is a member of the same large-base filter, so their finite
intersection is a member too. Restriction from $Q'\wedge R$
to each $Q'\wedge R_j$ changes no bundle or link. The preceding
construction applies, and forgetting $R$ in (DEX.17) gives the required
selector over $Q'$. Thus the common class and selector follow from the
literal eventual row and existence premises. Their actual arithmetic
suppliers, and recursive existence in this class, are still required.

\paragraph{DEX9: construct the family on all real indices from actual kits.}
\label{claim:DEX9}
Suppose $S$ is the supplied unbounded set of allowed scales, with
fixed $J$, and $x_1>1$. For every allowed scale $s$ with
$s.x\ge x_1$, suppose an original nonempty kit $K_s$ and its
\texttt{NumAt} records are supplied. Choose one allowed $y\ge x_1$.
For each real $x$, use scale $x$ if $x>x_1$ and $x\in S$;
otherwise use the fixed fallback scale $y$. Denote this scale by $s_x$.
Define
\[
 I_x=\{i\in K_{s_x}.I:\ x\le x_1\ \text{or}\ x\in S\}.    \tag{DEX.18}
\]
It is nonempty on $S$, and at $x=0$, as the family definition requires.
For late $x$ outside $S$ it is empty. Consequently every actual late
member has $s_x.x=x$; the fallback creates no false statement about
an unallowed late scale.

For $i\in I_x$, copy the point's $D,L,G_x,G_y$, extra record,
local objects and ties. Copy the common numerical constants from $C$.
The kit's equality clauses identify ranks, Artin maps, primes, root
discriminants, canonical infima, exponents, selected weights and
period data at this fixed $x$. Its word maps give the same word codes,
and its coverage gives the family coverage by retaining the proof in
(DEX.18). The chosen \texttt{NumAt} records give all parameter, size,
top-bound and constant identities. In particular $b_{\rm top}$ is
independent of the selected kit index; the original
$\mathrm{LaG}\ge x^\epsilon$ inequality gives the required lower
bound for it. These are exactly the fields of \texttt{famOfKits}.

Every supplied point property therefore holds at every actual member
by substitution, with no new transport of a lattice or metric.
The common lists are $F.c_{\rm pre}=C.c_{\rm pre}$,
$F.c_{\rm mid}=C.c_{\rm mid}$ and $F.c_{\rm last}=C.c_{\rm last}$.
For the output list its index 15 is $\max(C_{pde},J)$.
Any fallback kit supplies $J\le C_{pde}$ through its host-size
condition, so this is $C_{pde}$ and the output list also agrees.
Thus the family has the exact parameter lists needed for
\texttt{Large}, not merely comparable lists. The assertions
$F.\mathrm{Win}=S$ and $F.J=J$ follow directly from the definition.

\paragraph{DEX10: every field of the compatible \texttt{DesignChoices} join.}
\label{claim:DEX10}
Assume the literal \texttt{KitCPW} premise of \texttt{DesignS}, for
the point predicate
$((\mathrm{withDuty}(\mathrm{withObj}(Q_\kappa)))\wedge
\mathrm{WallPt})\wedge G.\mathrm{pt}$, and fix $\lambda>1$.
We prove its \texttt{DesignChoices} conclusion, using DEX3 for
\texttt{MinkDiv}. All kit and point conditions here are on the
same original objects.

Take $\epsilon_0$ to be the minimum of the kit bound and $1/53$.
For $0<\epsilon\le\epsilon_0$, the kit gives constants
$C_{T2},C_8,C_5,c_6,C_b$ and threshold $K_0$, before any family
is chosen. For the requested $K$, first replace it by
$\mathrm{thrTu}(K,c_T(\epsilon,9))$: this replaces entry 13 in the
argument lists of its middle and last bounds by $c_T(\epsilon,9)$.
It is the exact threshold transformation for the subsequent
\texttt{setTu} operation.

There is an elementary standard parameter record large for this
threshold and $K_0$. Explicitly choose its pre-parameter $p$ at least
$1000$ and both pre-bounds. Choose its middle parameter $a$ at least
$1$ and both middle bounds on the already fixed output list. Then
choose its last parameter $b$ at least $1$ and both last bounds on
that list followed by the middle list. In the source's
\texttt{stdConst}, every pre-entry is $p$, every middle-entry is $a$,
and every last-entry is $b$; its fixed entries are
$C_t=401$, $C_c=2$, $C_{YL}=9$, $c_w=9/10$, $e_P=12$,
$C_r=C_{dg}=C_{ccK}=1$, and the stated fixed gap constants.
Its other positive gap entries are $a,b,b^{-1}$, in their prescribed
positions. Substitution verifies \texttt{Std}, both largeness
conditions and nonnegativity of the common gap record. No later
choice enters an earlier bound.

For this $C$, take the maximum of the kit's $c_D$ threshold and
SZ9's numerical threshold, then the maximum of their $c_h$
thresholds, then $x_1=\max(2,x_{\rm kit},x_{\rm num})$.
Choose one supplied kit at each allowed late scale. SZ9 gives its
\texttt{NumAt}, \texttt{ChTT} and \texttt{ChCode}; DEX9 constructs
one family $F$. Its common choice is
\[
       \kappa=(\epsilon,C,c_D,c_h,C_5,c_6,C_b),
       \qquad C_H=C.P.C_F\ge1.                          \tag{DEX.19}
\]
The copied point conditions give the original class property and
$G.\mathrm{pt}$. The same set of windows supplies $G.\mathrm{cls}$.
The \texttt{NumAt} identities fix the page constants and windows,
in particular $C_{YL}=9$. They also give
\texttt{ThickSetEventually}, \texttt{ChosenTT},
$\texttt{ChosenBad}(C_b)$, $\texttt{ChosenQD}(C_5)$ and
$\texttt{ChosenCode}(3,C_H)$. Their translations to the family
variables use $L.sc.x=x$ at late indices and $L.sc.\epsilon=\epsilon$;
the threshold $F.x_1+1$ makes these substitutions legitimate.

The \texttt{WallPt} and original duty clauses give
\texttt{BoxesF} and \texttt{QDRowsF}, as well as the chosen wall
and twist properties, by their definitions. The \texttt{MinkPt}
clause, applied to the precise rational divisor of DEX3, gives
\texttt{DGamDvdF}; its tame clause gives \texttt{TameDistinctF}.
The middle-list equality identifies $F.C_F=C_H$. The word count
is literally
\[
                 |\mathrm{Word}|=\lceil3\Gamma^2x^{C_H}\rceil,
\]
so the ceiling inequality yields
$|\mathrm{Word}|\le3\Gamma^2x^{C_H}+1$, proving
\texttt{CodeUpF} with its original additive one.

Finally set $F'=\mathrm{setTu}(F,c_T(\epsilon,9))$.
This changes only the scalar field $C_{Tu}$ of the family record.
All the actual data and point predicates just verified are unchanged.
The definition of \texttt{thrTu} gives exactly
$F'.\mathrm{Large}(K,C_H)$, and $C_{YL}=9$ gives
$c_T(\epsilon,F'.C_{YL})\le F'.C_{Tu}$ with equality.
For the source $\texttt{Link}(\kappa,F')$, its output clause first
sets $C_{Tu}$ back to $C.CTu$; hence the earlier exact output-list
identity verifies that clause as well. This proves all fields of
the literal \texttt{DesignChoices} conclusion, on the original
target-compatible class and with one witness family.

\paragraph{DEX11: the recursive input that still produces the compatible class.}
\label{claim:DEX11}
Let $Q_O(\kappa,D,O)$ be the original object predicate chosen in
DEX8, or any supplied predicate of that type. Assume its actual
$\mathrm{RecExP}(\mathrm{withQD}(Q_O))$ and an actual
$\mathrm{HostExW}(\lambda,G.\mathrm{req})$, with $\lambda>1$.
The reconstructed RF10 point join applies: AC supplies its
Harish--Chandra and level constructors, SZ its initial sizes, QR25
the required level-prime assertion, and the assumed recursive
object supplies the actual $D,Y$ after every level choice. It gives
\texttt{ObjCP} for precisely $Q_O$, including the unchanged
\texttt{withQD} object and all original host, code and canonical
instance equalities.

Here the quantifier ``after every level choice'' is essential:
the positive integers $d_u$, positive $D_*$ and logarithmic window
are selected before the universal $Lv$. The recursive premise
then works for each $Lv$ with those same values. RF10 can therefore
use the actual prime levels obtained from QR25. An existential
recursive datum for one unrelated level would not suffice.

Forget $Q_O$ temporarily. HX1 supplies the required small prime,
and HW1--HW2, using GW5, construct the degree-selected host from
this point constructor and the supplied \texttt{HostExW}.
The injective word code can also be chosen uniformly, before the
host and all its positions. Put $H=m^5f x^\epsilon$ and
$M=\lceil3\exp(2C_\Gamma H)x^{C_F}\rceil_+$. The literal rank
windows give
\[
 \begin{gathered}
 m\le2x^{1/2-\epsilon/16},\qquad
 f\le4x^{2+8\epsilon+\epsilon/100},\\
 H\le128x^{9/2+9\epsilon},\qquad
 k\ge x^{9/2+10\epsilon}.
 \end{gathered}
                                                               \tag{DEX.21}
\]
For fixed positive $\epsilon$ and nonnegative $C_\Gamma,C_F$,
the ceiling estimate $M<2(3\exp(2C_\Gamma H)x^{C_F})$ gives
\[
 \log M\le\log6+256C_\Gamma x^{9/2+9\epsilon}
                         +C_F\log x=o(x^{9/2+10\epsilon}).
                                                               \tag{DEX.22}
\]
Choose one threshold depending only on these three parameters so
that $k\ge6$ and the right side is at most $k/4$. Since
$\log2>1/2$, we have $k/4\le(k-3)\log2$, hence
$0<M\le2^{k-3}\le2^{k-2}$. Enumerating the Boolean cube therefore
gives an injection $\operatorname{Fin}M\hookrightarrow
\{0,1\}^{\operatorname{Fin}(k-2)}$. This threshold is fixed before
the scale, the host and every $u_0$; it supplies the literal
\texttt{exists\_injective\_words} input of \texttt{kitS}.

This also corrects the sharper intermediate estimate in the historical
RF.45--RF.46: the rank window for $f$ contains $\epsilon/100$,
which those displayed exponents omitted. Use (DEX.21) there, with
positive exponent gap $\epsilon$, and choose
$x^\epsilon\ge2048C_\Gamma$ for its first eighth-bound.
The logarithmic eighth-bound is unchanged. The argument above still
proves the stated RF.48 conclusion $\#\mathrm{Word}\le2^{k-3}$.
Only that historical intermediate estimate is superseded; the
word bound and its uniform quantifier order are preserved.

Fix that host, its degree pair and the resulting injective word code, then
invoke the original point constructor once for each $u_0<n$.
The common host and code identify their word values; the source
canonical infimum and filling formulas give the exact kit equalities
as proved in HW2. The point indexed by $u_0$ supplies kit coverage
at that position. HW1 gives the wall condition for every point,
and the copied target-host property gives $G.\mathrm{pt}$.
Combining thresholds in the order (DEX.16) yields the literal
\texttt{KitCPW} used in DEX10. That claim gives
\texttt{DesignChoices} for
\[
 \mathrm{WithTw}\{F:\ \mathrm{clsP}(Q_O)(F)
                                  \wedge G.\mathrm{cls}(\lambda,F)\}.
                                                               \tag{DEX.20}
\]
This is the full conditional \texttt{recS} route, with the original
object predicate and target class, not just a nonempty family.

The remaining recursive supply must be read against the latest
chapters. QW9 already proves its actual logarithmic window, and
MF20 discharges the mass-ratio premise on the same data. Therefore
IG14's historical window condition is closed. What remains in
that join is the actual \texttt{RecAtG} core and the eventual
$Q_O$ property, with their original objects and dependent equalities.
The immediate recursive output has predicate
$\kappa.\mathrm{Pos}\Rightarrow Q_O$; for the particular
$Q_O=\mathrm{QOf}$ of DEX8 this is equivalent to $Q_O$, since
its definition already begins with that same positivity guard.
Thus this simplification neither discards a hypothesis nor strengthens
the recursive producer.
DEX8 supplies the finite conjunction and selection step once its
bundle-existence and row premises are proved. DEX7 then supplies
\texttt{DesignExists} once the same class also has the required
\texttt{TuSize} implication and actual rank formula. These are
the residual construction obligations; DEX1--DEX11 do not assume
them away or identify independently chosen bundles.

\end{document}
